% 第 17 章 Pattern Library Manager
\bichapter{图案库管理器}{Pattern Library Manager}
\label{chap:plm}

本章介绍 IC Validator 图案库管理器（Pattern Library Manager）。

IC Validator 图案库管理器按下列各节说明：
\begin{itemize}
  \item 入门
  \item 创建新图案库
  \item 编辑图案
  \item 导出与导入图案库
  \item 在图案库管理器中运行图案匹配
  \item 限制与局限
\end{itemize}

% ============================================================
\section{入门}
\subsection*{Getting Started}

IC Validator 图案库管理器（Pattern Library Manager，PLM）是图形用户界面（GUI），提供另一种方式为 IC Validator 图案匹配运行构建图案库，而无需使用批处理模式通过 \cmd{pattern\_learn()} 函数生成图案库。

\subsection{支持的 IC Validator WorkBench 版本}
\subsubsection*{Supported Versions of IC Validator WorkBench}

图案库管理器通过套接字或端口连接与 IC Validator WorkBench 工具通信。

IC Validator N-2017.12-SP2 及更高版本支持 IC Validator WorkBench J-2014.06 及更高版本。

\subsection{启动图案库管理器}
\subsubsection*{Invoking the Pattern Library Manager}

确保已正确设置 \cmd{ICV\_INSTALL\_DIRECTORY} 安装目录。

可在 shell 命令行或 IC Validator WorkBench 版图编辑器中启动图案库管理器。

\begin{enumerate}
  \item 命令行语法为：
\begin{lstlisting}[style=shell]
icv_plm
\end{lstlisting}

\begin{noteBox}
若分别从命令行启动图案库管理器与 IC Validator WorkBench，二者必须位于同一目录路径（同一台机器上），才能建立连接。关于如何将图案库管理器连接到 IC Validator WorkBench，见「与 IC Validator WorkBench 通信」一节。
\end{noteBox}

  \item 要从 IC Validator WorkBench 访问图案库管理器，在命令窗口中 source IC Validator 的 \cmd{plm\_button.tcl} 文件：
\begin{lstlisting}[style=shell]
% source $env(ICV_INSTALL_DIRECTORY)/contrib/plm_button.tcl
\end{lstlisting}

  工具栏会添加 \textbf{ConnectPlm} 按钮。单击该按钮即可启动图案库管理器并建立连接，如图~\ref{fig:plm-connect} 所示。
\end{enumerate}

\figplaceholder{Figure 111: ConnectPlm Button}{ConnectPlm 按钮}{fig:plm-connect}

\begin{noteBox}
也可将该 \cmd{plm\_button.tcl} 文件放入 IC Validator WorkBench 启动文件，以免每次都重新设置连接。
\end{noteBox}

设置 \cmd{ICVWB\_USER} 环境变量后，IC Validator WorkBench 工具会自动连接到图案库管理器。

对于 IC Validator WorkBench M-2016.03-SP2 及更新版本：
\begin{lstlisting}[style=shell]
setenv ICVWB_USER ICV_PLM
\end{lstlisting}

工具栏会添加一个黄色按钮。单击该按钮即可启动图案库管理器并建立连接。

\figplaceholder{Figure 112: Pattern Library Manager Toolbar}{图案库管理器工具栏}{fig:plm-toolbar}

\subsection{与 IC Validator WorkBench 通信}
\subsubsection*{Communicating With IC Validator WorkBench}

若图案库管理器已启动，要与 IC Validator WorkBench 建立连接：

\begin{enumerate}
  \item 在 IC Validator WorkBench 命令窗口中输入下列命令：
\begin{lstlisting}[style=shell]
user_socket open 0
\end{lstlisting}

  该命令返回一个数字。这是 IC Validator WorkBench 建立的打开套接字。

\figplaceholder{Figure 113: Open Socket Connection}{打开套接字连接}{fig:plm-socket}

  \item 在图案库管理器中，选择 \textbf{Tools > Communication}。输入端口号并单击 \textbf{OK}。

\figplaceholder{Figure 114: Communication Window}{通信窗口}{fig:plm-comm}

  \item 此时图案库管理器已连接到 IC Validator WorkBench。
\end{enumerate}

\subsection{图案库管理器 GUI}
\subsubsection*{Pattern Library Manager GUI}

主 GUI 如图所示。表~\ref{tab:plm-menus} 给出该 GUI 四个下拉菜单的主要功能。

\begin{longtable}{@{}p{0.22\textwidth} p{0.70\textwidth}@{}}
\caption{图案库管理器的下拉菜单}\label{tab:plm-menus}\\
\toprule
\textbf{菜单} & \textbf{说明} \\
\midrule
\endfirsthead
\caption[]{图案库管理器的下拉菜单（续）}\\
\toprule
\textbf{菜单} & \textbf{说明} \\
\midrule
\endhead
\bottomrule
\endfoot
File
  & 创建、打开或保存图案库。导入 Pattern Database 格式的图案库，并导出 Pattern Library Manager 格式的图案库 \\
Edit
  & 在图案库中创建、添加或编辑图案 \\
Tool
  & 与 IC Validator WorkBench 建立通信，并在图案库管理器中执行图案匹配 \\
Help
  & 显示关于界面语言与用户指南的信息 \\
\end{longtable}

% ============================================================
\section{创建新图案库}
\subsection*{Creating a New Pattern Library}

\textbf{File} 菜单允许创建用于添加新图案的新图案库。表~\ref{tab:plm-file-menu} 汇总可用选项。

\begin{longtable}{@{}p{0.22\textwidth} p{0.70\textwidth}@{}}
\caption{File 菜单命令}\label{tab:plm-file-menu}\\
\toprule
\textbf{选项} & \textbf{说明} \\
\midrule
\endfirsthead
\caption[]{File 菜单命令（续）}\\
\toprule
\textbf{选项} & \textbf{说明} \\
\midrule
\endhead
\bottomrule
\endfoot
New
  & 创建新图案库 \\
Open
  & 打开既有图案库 \\
Save
  & 保存当前库 \\
Export
  & 将图案库输出为 Pattern Database 格式 \\
Import
  & 将 Pattern Database 格式的库读入图案库管理器 \\
Close
  & 关闭当前库 \\
Exit
  & 退出图案库管理器 \\
\end{longtable}

\subsection{创建新图案库名称}
\subsubsection*{Creating a New Pattern Library Name}

在向新图案库添加新图案之前，必须先创建图案库名称。

\begin{enumerate}
  \item 选择 \textbf{File > New}。出现对话框。以默认名称 \texttt{New1} 创建新库名。
  \item 必要时更改库名。
  \item 指定保存图案库的路径。
  \item 填写所需的图案库属性，例如图案模糊度（pattern fuzziness）、图案层数量，以及均匀模糊尺寸（uniform fuzzy size，仅用于 \cmd{PM\_EDGE\_UNIFORM} 模式）。
  \item 单击 \textbf{More} 打开高级设置，例如图案类型与可选图案标记层数量。
  \item 单击 \textbf{OK} 创建新图案库（名称）。IC WorkBench EV Plus 窗口中会创建对应的空版图（以图案库名称命名），并含多层。图案库管理器主窗口中会出现列出已定义图案库参数的摘要表，如图~\ref{fig:plm-new-lib} 所示。
  \item 现在可以向新图案库添加图案。添加新图案时，每个图案都会出现在已创建的图案库下。
\end{enumerate}

\figplaceholder{Figure 115: Creating a New Pattern Library}{创建新图案库}{fig:plm-new-lib}

\subsection{添加新图案}
\subsubsection*{Adding New Patterns}

可使用 \textbf{Edit} 菜单中的三个命令向图案库添加新图案，如表~\ref{tab:plm-edit-menu} 所示。

\begin{longtable}{@{}p{0.32\textwidth} p{0.60\textwidth}@{}}
\caption{Edit 菜单命令}\label{tab:plm-edit-menu}\\
\toprule
\textbf{命令} & \textbf{说明} \\
\midrule
\endfirsthead
\caption[]{Edit 菜单命令（续）}\\
\toprule
\textbf{命令} & \textbf{说明} \\
\midrule
\endhead
\bottomrule
\endfoot
New Pattern
  & 创建图案版图并将其添加为新图案 \\
New Pattern By Clip
  & 通过版图上的裁剪窗口生成图案 \\
New Pattern By Layer
  & 从版图选定区域提取多个图案 \\
Copy
  & 从图案库复制图案 \\
Paste
  & 将复制的图案粘贴到另一图案库 \\
Rename
  & 重命名图案 \\
Delete
  & 从图案库删除图案 \\
\end{longtable}

\subsubsection{通过裁剪捕获新图案}
\paragraph*{Capturing a New Pattern by Clipping}

可从已定义图案层与图案标记的版图创建新图案。图案边界定义为裁剪窗口。

要捕获新图案，从包含多边形的版图开始。

\begin{enumerate}
  \item 在图案库管理器中，打开既有图案库或创建新图案库条目。
  \item 在 IC WorkBench EV Plus 中打开版图。
  \item 浏览到图案位置，并缩放到合适显示，使所有待提取多边形显示在显示窗口中。
  \item 如图~\ref{fig:plm-clip} 所示，在图案库管理器中选择 \textbf{Edit > New Pattern By Clip}。出现对话框。为该图案输入名称，并指定图案标记层与图案层的层号与数据类型。单击 \textbf{OK}。

\figplaceholder{Figure 116: Capturing a New Pattern By Clipping}{通过裁剪捕获新图案}{fig:plm-clip}

\begin{noteBox}
默认情况下，图案范围外的多边形不会被捕获到生成的图案中。若要在图案最终确定前保留图案范围外的周围区域以便进一步修改，请在单击 \textbf{OK} 之前选择 \textbf{Keep Context}。
\end{noteBox}

  \item 在 IC WorkBench EV Plus 中绘制裁剪窗口以定义图案范围。出现新版图窗口。每条多边形边从 0 开始标注编号。边标签作为文本层创建在层列表下。新图案被添加到图案库，如图~\ref{fig:plm-clip-window} 所示。

\figplaceholder{Figure 117: Drawing a Clipping Window}{绘制裁剪窗口}{fig:plm-clip-window}

\begin{noteBox}
在步骤~4 中单击 \textbf{OK} 之后，IC WorkBench EV Plus 中唯一允许的操作是绘制裁剪窗口以定义图案范围。
\end{noteBox}

  \item 如需要，在 \textbf{Pattern Property} 选项卡中定义图案属性。更多信息见「向图案附着图案属性」一节。
  \item 如需要，在 \textbf{Tolerance} 选项卡中添加边约束。更多信息见「向图案添加边容差」一节。
  \item 若同一版图中有多个位置，重复前述步骤以生成更多新图案。完成后，在图案库管理器中选择 \textbf{File > Save}，将更改保存到图案库文件。
\end{enumerate}

由此生成包含来自版图的一个或多个图案的图案库。

\subsubsection{在版图中提取图案}
\paragraph*{Extracting Patterns in a Layout}

可从已定义图案层、图案标记与图案范围层的版图创建多个图案。

\begin{enumerate}
  \item 在图案库管理器中，打开既有图案库或创建新图案库条目。
  \item 在 IC WorkBench EV Plus 中打开版图。
  \item 浏览到版图位置，并缩放到合适显示，使版图的正确部分显示在显示窗口中。
  \item 如图~\ref{fig:plm-extract} 所示，在图案库管理器中选择 \textbf{Edit > New Pattern By Layer}。出现对话框。输入图案前缀名称，并指定图案范围、图案标记与图案层的层号与数据类型。若已定义可选忽略区域，也可指定。单击 \textbf{OK}。

\figplaceholder{Figure 118: Extracting Patterns in a Layout}{在版图中提取图案}{fig:plm-extract}

\begin{noteBox}
默认情况下，图案范围外的多边形不包含在生成的图案中。若要保留图案范围外的周围区域，请在单击 \textbf{OK} 之前选择 \textbf{Keep Context}。
\end{noteBox}

  \item 在 IC WorkBench EV Plus 中绘制裁剪窗口，以选择用于生成新图案的版图部分。图案被提取并列在当前打开或选定的图案库下，如图~\ref{fig:plm-select-area} 所示。

\figplaceholder{Figure 119: Selecting Part of Layout for Generating New Patterns}{选择版图部分以生成新图案}{fig:plm-select-area}

\begin{noteBox}
在步骤~4 中单击 \textbf{OK} 之后，IC WorkBench EV Plus 中唯一允许的操作是绘制裁剪窗口，以定义用于图案提取的版图区域。
\end{noteBox}

  \item 如需要，在 \textbf{Pattern Property} 选项卡中定义图案属性。更多信息见「向图案附着图案属性」一节。
  \item 如需要，在 \textbf{Tolerance} 选项卡中添加边约束。完整信息见「向图案添加边容差」一节。
  \item 完成后，在图案库管理器中选择 \textbf{File > Save}，将更改保存到图案库文件。
\end{enumerate}

由此生成包含来自版图的多个图案的图案库，如图~\ref{fig:plm-multi-patterns} 所示。

\figplaceholder{Figure 120: Pattern Library With Multiple Patterns}{含多个图案的图案库}{fig:plm-multi-patterns}

\subsubsection{从零生成新图案}
\paragraph*{Generating a New Pattern From Scratch}

当没有既有版图可供图案提取时，可在 ICWBEV 中绘制图案版图以创建新图案。

\begin{enumerate}
  \item 在图案库管理器中，打开既有图案库或创建新图案库条目。
  \item 在图案库管理器中，选择 \textbf{Edit > New Pattern}。出现对话框。输入待生成图案的名称，然后单击 \textbf{OK}，如图~\ref{fig:plm-new-scratch} 所示。

\figplaceholder{Figure 121: Generating a New Pattern From Scratch}{从零生成新图案}{fig:plm-new-scratch}

  \item 出现警告消息。这是因为 IC WorkBench EV Plus 中尚无与新创建图案名称关联的图形。单击 \textbf{OK} 关闭该消息。

\figplaceholder{Figure 122: No Graphics Warning Message}{无图形警告消息}{fig:plm-no-graphics}

  \item 在 IC WorkBench EV Plus 中创建图案版图，如图~\ref{fig:plm-create-layout} 所示。

\figplaceholder{Figure 123: Creating the Pattern Layout}{创建图案版图}{fig:plm-create-layout}

  \item 关于版图创建与编辑的更多信息，见 Help 菜单中的 \textit{IC WorkBench EV Plus User Guide}。见图~\ref{fig:plm-help-menu}。

\figplaceholder{Figure 124: IC WorkBench EV Plus User Guide Help Menu}{IC WorkBench EV Plus 用户指南 Help 菜单}{fig:plm-help-menu}

  \item 在 IC WorkBench EV Plus 中单击 \textbf{PlmRefresh}，用为该图案创建的全部层更新图案库。图案库管理器中出现对话框，如图~\ref{fig:plm-update-dlg} 所示。单击 \textbf{Yes}。

\figplaceholder{Figure 125: Dialog Box for Updating Pattern Library}{更新图案库对话框}{fig:plm-update-dlg}

  图案被添加到图案库，如图~\ref{fig:plm-pattern-added} 所示。

\figplaceholder{Figure 126: Pattern Added to Pattern Library}{图案已添加到图案库}{fig:plm-pattern-added}

  \item 如需要，在 \textbf{Pattern Property} 选项卡中定义图案属性。更多信息见「向图案附着图案属性」一节。
  \item 如需要，在 \textbf{Tolerance} 选项卡中添加边约束。更多信息见「向图案添加边容差」一节。
  \item 重复前述步骤，在同一图案库文件下生成更多图案。完成后，在图案库管理器中选择 \textbf{File > Save}，将更改保存到图案库文件。
\end{enumerate}

% ============================================================
\section{编辑图案}
\subsection*{Editing a Pattern}

可在图案库管理器内部编辑图案：向该图案附着一个或多个文本属性，或为非均匀匹配模式指定边容差。

也可在 IC WorkBench EV Plus 中编辑图案版图：添加、删除或移动多边形、边与顶点。

\subsection{向图案添加边容差}
\subsubsection*{Adding an Edge Tolerance to a Pattern}

图案库管理器允许为非均匀模糊匹配模式指定边容差（约束）。约束使您能够将图案的不同变体匹配到已定义图案。每条边不应有超过一个边约束。

要指定边容差：

\begin{enumerate}
  \item 打开包含待编辑图案的图案库。
  \item 单击图案名称以选择待编辑图案。图案库管理器中出现 \textbf{Tolerance} 与 \textbf{Pattern Property} 选项卡，图案版图显示在 IC WorkBench EV Plus 中。
  \item 从下拉框中选择要指定容差的边。若所选边为水平边，则该边被限制为垂直移动；若所选边为竖直边，则该边被限制为水平移动。

\begin{noteBox}
可使用 IC WorkBench EV Plus 中显示的边编号作为参考，以识别边编号。
\end{noteBox}

  \item 单击增大或减小按钮，以纳米为单位指定内侧与外侧限值。若容差区间包含该值，选择 \textbf{Inclusive}，如图~\ref{fig:plm-edge-tol} 所示。

\figplaceholder{Figure 127: Adding an Edge Tolerance}{添加边容差}{fig:plm-edge-tol}

  \item 单击 \textbf{Add}，将容差添加到图案库。边容差出现在 IC WorkBench EV Plus 的层 \texttt{(9001:0)} 上。
  \item 选择 \textbf{File > Save}，将更改保存到图案库。
  \item 单击 \textbf{Update} 修改既有边容差，单击 \textbf{Delete} 删除边容差。
  \item 也可在 \textbf{Edge2Edge} 窗格中通过选择两条边来指定边到边容差，如图~\ref{fig:plm-e2e-tol} 所示。

\figplaceholder{Figure 128: Adding an Edge-to-Edge Tolerance}{添加边到边容差}{fig:plm-e2e-tol}

\begin{noteBox}
要使边到边约束被识别，所选两条边中至少有一条必须已定义边容差。
\end{noteBox}
\end{enumerate}

\subsection{向图案附着图案属性}
\subsubsection*{Attaching Pattern Properties to a Pattern}

图案库管理器允许在图案库创建期间向图案附着属性。找到匹配时，这些属性会传递到输出图案标记。这使您能够使用其他 IC Validator runset 函数，基于属性名与属性值对输出图案标记进行后处理。

要向图案附着属性：

\begin{enumerate}
  \item 打开包含待编辑图案的图案库。
  \item 单击图案名称以选择待编辑图案。出现 \textbf{Tolerance} 与 \textbf{Pattern Property} 选项卡。单击 \textbf{Pattern Property} 选项卡以激活，如图~\ref{fig:plm-prop-tab} 所示。

\figplaceholder{Figure 129: Pattern Property Tab}{Pattern Property 选项卡}{fig:plm-prop-tab}

  \item 若属性为字符串类型，将属性名与值添加到 \textbf{String Property} 字段；若为双精度类型，则添加到 \textbf{Double Property} 字段。对于图~\ref{fig:plm-attach-prop} 中的示例，\texttt{hs\_type} 属性为字符串类型，赋值为 \texttt{space}；\texttt{min\_cd} 属性为双精度类型，赋值为 \texttt{0.056}。

\figplaceholder{Figure 130: Attaching Pattern Properties to a Pattern}{向图案附着图案属性}{fig:plm-attach-prop}

  \item 单击 \textbf{Add}，将属性添加到图案库。
  \item 选择 \textbf{File > Save}，将更改保存到图案库。
  \item 单击 \textbf{Update} 修改既有属性，或单击 \textbf{Delete} 删除属性。
\end{enumerate}

\subsection{在 IC WorkBench EV Plus 中编辑图案版图}
\subsubsection*{Editing a Pattern Layout in IC WorkBench EV Plus}

可在 IC WorkBench EV Plus 中修改图案版图：为图案层、图案范围、图案标记与忽略区域层添加、移动或删除多边形与形状。关于如何编辑版图的更多细节，见 Help 菜单中的 \textit{IC WorkBench EV Plus User Guide}。

移除细条（sliver）的示例：

\begin{enumerate}
  \item 打开包含待编辑图案的图案库。
  \item 单击图案名称以选择待编辑图案。图案 cell 被选中并出现在 IC WorkBench EV Plus 中。
  \item 相应编辑该图案。例如，若要移除图~\ref{fig:plm-edit-layout} 中 \texttt{pat\_2} 的细条，执行下列步骤：

\figplaceholder{Figure 131: Editing a Pattern Layout}{编辑图案版图}{fig:plm-edit-layout}

  \begin{enumerate}[label=\alph*.]
    \item 要选择形状，选择 \textbf{Tool > Select}，或单击 Tools 工具栏上的 \textbf{Select} 按钮并单击该形状。
    \item 要从 \texttt{pat\_2} 移除所选形状，选择 \textbf{Edit > Delete}，或按键盘上的 Delete 键。
    \item 在 IC WorkBench EV Plus 中单击 \textbf{PlmRefresh}，以更新对 \texttt{pat\_2} 所做的更改。出现消息，如图~\ref{fig:plm-refresh} 所示。单击 \textbf{Yes}。
  \end{enumerate}

\figplaceholder{Figure 132: Pattern Library Manager Refresh Button}{图案库管理器刷新按钮}{fig:plm-refresh}

  \item 在图案库管理器中，选择 \textbf{File > Save} 保存图案库，如图~\ref{fig:plm-save} 所示。

\figplaceholder{Figure 133: Saving the Pattern Library}{保存图案库}{fig:plm-save}
\end{enumerate}

% ============================================================
\section{导出与导入图案库}
\subsection*{Exporting and Importing a Pattern Library}

在 IC Validator 工具中，可使用基于 GUI 的图案库管理器创建图案库；也可通过运行 IC Validator 作业、使用 \cmd{pattern\_learn()} 函数创建图案库。

在图案库管理器中生成的图案库由两部分组成：用于存储图形信息的 OASIS 目录，以及用于存储边容差与边到边约束的 XML 格式文本文件。

通过运行 IC Validator 作业、使用 \cmd{pattern\_learn()} 函数创建的图案库由三部分组成：名为 \cmd{pattern.dat} 的二进制文件、名为 \cmd{log} 的文本文件（记录图案库创建参数与图案统计），以及在图案匹配模式设为 \cmd{PM\_EDGE\_NONUNIFORM} 时可选的、用于存储边容差与边到边约束的 XML 格式文本文件。

\cmd{pattern\_match()} 函数在 IC Validator 运行中执行图案匹配。图案匹配运行前图案库必须就绪，且必须为二进制格式，也称为图案数据库（Pattern Database，PDB）格式。

表~\ref{tab:plm-formats} 给出这两种图案库格式的摘要比较。

\begin{longtable}{@{}p{0.28\textwidth} p{0.38\textwidth} p{0.26\textwidth}@{}}
\caption{图案库格式摘要比较}\label{tab:plm-formats}\\
\toprule
\textbf{模式} & \textbf{格式} & \textbf{导入/导出} \\
\midrule
\endfirsthead
\caption[]{图案库格式摘要比较（续）}\\
\toprule
\textbf{模式} & \textbf{格式} & \textbf{导入/导出} \\
\midrule
\endhead
\bottomrule
\endfoot
使用 \cmd{pattern\_learn()} 函数创建的图案库
  & Pattern Database 格式包括：
    (1)~\cmd{pattern.dat}（二进制）；
    (2)~\cmd{log}；
    (3)~可选 XML 文件，在匹配模式为 \cmd{PM\_EDGE\_NONUNIFORM} 时存储边容差
  & 必须导入图案库管理器才能编辑 \\
在图案库管理器中创建的图案库
  & Pattern Library Manager 格式包括：
    (1)~OASIS，用于存储图形信息；
    (2)~XML 文件，用于存储边与边到边容差
  & 必须从图案库管理器导出，才能在 IC Validator 图案匹配运行中使用该图案库 \\
\end{longtable}

\subsection{导入 Pattern Database 格式}
\subsubsection*{Importing a Pattern Database Format}

Pattern Database 格式的图案库必须导入后，才能使用图案库管理器进行编辑。

\begin{enumerate}
  \item 在图案库管理器中，选择 \textbf{File > Import}。
  \item 浏览到 Pattern Database 格式图案库的位置。浏览到将转换为 Pattern Library Manager 格式的图案库的保存位置，并为其输入名称，如图~\ref{fig:plm-import} 所示。

\figplaceholder{Figure 134: Importing a Pattern Library}{导入图案库}{fig:plm-import}

\begin{noteBox}
若 Pattern Database 格式的图案库已锁定（加密），导入将失败。
\end{noteBox}

  \item 单击 \textbf{OK}。图案库被转换为 Pattern Library Manager 格式并加载到图案库管理器中。
\end{enumerate}

\subsection{导出 Pattern Library Manager 格式}
\subsubsection*{Exporting a Pattern Library Manager Format}

Pattern Library Manager 格式的图案库必须导出后，才能在 IC Validator 图案匹配运行中使用。

\begin{enumerate}
  \item 在图案库管理器中，选择 \textbf{File > Open} 加载待导出的图案库，或单击选择已在图案库管理器中打开的图案库。
  \item 选择 \textbf{File > Export}。
  \item 浏览到将转换为 Pattern Database 格式的图案库的保存位置，并为其输入名称，如图~\ref{fig:plm-export} 所示。

  若不同取向的图案也应视为匹配，选择 \textbf{pattern reflect} 与 \textbf{pattern rotate}。

  若在 \cmd{PM\_EXACT} 匹配模式下允许额外多边形，选择 \textbf{ignore extra polygons}。

\figplaceholder{Figure 135: Exporting a Pattern Library}{导出图案库}{fig:plm-export}

  \item 单击 \textbf{OK}。图案库被转换为 Pattern Database 格式，并按步骤~3 中定义的名称与路径保存到磁盘。出现 Export Result 窗口，如图~\ref{fig:plm-export-result} 所示。单击 \textbf{OK} 关闭该窗口并完成导出过程。

\figplaceholder{Figure 136: Pattern Library Export Result}{图案库导出结果}{fig:plm-export-result}
\end{enumerate}

该图案库已可在 IC Validator 运行中使用。

% ============================================================
\section{在图案库管理器中运行图案匹配}
\subsection*{Running Pattern Matching in the Pattern Library Manager}

IC Validator 工具使用 \cmd{pattern\_match()} 函数在 IC Validator 运行中执行图案匹配。这意味着必须创建 IC Validator runset。关于如何创建 IC Validator 规则文件并在批处理模式下运行图案匹配的更多信息，见本用户指南「图案匹配」章。

除批处理模式外，图案库管理器还提供便捷方式：无需离开 GUI，即可使用已生成的图案库对设计版图运行图案匹配。

\begin{enumerate}
  \item 在图案库管理器中，对选定的图案库单击 \textbf{Pattern Match}，如图~\ref{fig:plm-pm-btn} 所示。

\figplaceholder{Figure 137: Performing Pattern Matching on the Selected Pattern Library}{对选定图案库执行图案匹配}{fig:plm-pm-btn}

  \item 弹出 Pattern Match 窗口，如图~\ref{fig:plm-pm-win} 所示。

\figplaceholder{Figure 138: Pattern Match Window}{Pattern Match 窗口}{fig:plm-pm-win}

  \begin{enumerate}[label=\alph*.]
    \item 浏览到输入版图的位置。选择设计格式。输入顶层 cell 名称。指定图案层。
    \item 可在 \textbf{Run Options} 列中指定 IC Validator 命令行选项。关于命令行选项的更多信息，见本用户指南「IC Validator 基础」章。
    \item 指定输出目录路径。
    \item 单击 \textbf{More} 显示附加选项。若希望按图案名对结果分组，选中 \cmd{group\_errors\_by\_pattern} 复选框。
    \item 单击 \textbf{OK}，对指定版图执行图案匹配。
  \end{enumerate}

  \item 图案匹配结果出现在对话框中，如图~\ref{fig:plm-pm-result} 所示。详细信息保存在输出目录中。单击 \textbf{OK} 关闭该对话框。

\figplaceholder{Figure 139: Pattern Matching Results Dialog Box}{图案匹配结果对话框}{fig:plm-pm-result}
\end{enumerate}

% ============================================================
\section{限制与局限}
\subsection*{Restrictions and Limitations}

下列功能目前在图案库管理器中不可用，但由批处理模式下的 \cmd{pattern\_learn()} 函数支持：
\begin{itemize}
  \item 指定由 IC Validator 工具根据用户定义的 ambit 值与 \cmd{ambit\_mode} 生成图案范围。
  \item 将边容差指定为用户定义输入层。
  \item 指定可选区域层。
  \item 指定自动锚定（auto anchoring）。
\end{itemize}
